\section{Designing Custom Thiasos}

A Thiasos is a fragment of a Patron's attention made manifest. It is a companion, a witness, and—for a Runic Warrior—often a weapon or shield. While this expansion provides dozens of pre-built Thiasos, the covenant between Runekeeper and Patron is personal. Sometimes, the standard options do not fit.

This chapter provides guidelines for Game Masters and players to design custom Thiasos, whether living, sentient, or unified with the Codex.

\subsection{The Three Questions}

Before building a custom Thiasos, answer three questions:

\subsubsection{What form does it take?}

\begin{itemize}
\item \textbf{Living:} A creature—wolf, horse, bird, serpent, hound. It breathes, eats, and can be wounded.
\item \textbf{Sentient:} An object—sword, shield, staff, bell, mirror. It speaks, remembers, but cannot act independently.
\item \textbf{Unified:} Both Codex and Thiasos. Rites are engraved on it. It serves as both record and witness.
\end{itemize}

\subsubsection{What Patron does it serve?}

The Thiasos's abilities should reflect its Patron's domain. A Thiasos of Grimmir might be a wolf-ancestor; a Thiasos of the Clockwork Monad might be a ticking brass compass. Refer to the Patron chapters for domain-specific guidance.

\subsubsection{What role does it play in the covenant?}

\begin{itemize}
\item \textbf{Combat:} Fights beside the Runekeeper. Has its own Harm track and actions.
\item \textbf{Utility:} Provides non-combat benefits—tracking, sensing, carrying messages.
\item \textbf{Witness:} Counts as a witness for Rites. May have perfect recall.
\item \textbf{Counsel:} Offers advice, warnings, or exposition. May have opinions the Runekeeper does not share.
\end{itemize}

\subsection{Designing a Living Thiasos}

\subsubsection{Core Statistics}

A living Thiasos has the following base statistics. Adjust based on Cap (2-5).

\begin{itemize}
\item \textbf{Body:} Cap + 1 (minimum 2)
\item \textbf{Wits:} Cap - 1 (minimum 1)
\item \textbf{Spirit:} Cap (minimum 1)
\item \textbf{Presence:} Cap - 2 (minimum 0)
\item \textbf{Harm:} Cap + 2 (minimum 4)
\item \textbf{Traits:} Two to three special abilities (see below)
\end{itemize}

\subsubsection{Sample Traits}

\begin{itemize}
\item \textbf{Tracker:} +2 dice to track a target by scent.
\item \textbf{Spirit Sense:} Can detect invisible or incorporeal creatures within Near range.
\item \textbf{Pack Hunter:} When the Runekeeper attacks a target, the Thiasos gains +1 die to attack the same target.
\item \textbf{Flight:} Can fly. +1 die to Stealth in natural terrain.
\item \textbf{Amphibious:} Can breathe underwater. No penalty for underwater combat.
\item \textbf{Regeneration:} Regains 1 Harm per day (faster than natural healing).
\item \textbf{Terrifying Presence:} Enemies within Near range test Resolve (DV 3) or become Shaken.
\item \textbf{Silent Movement:} Makes no sound. Cannot be detected by hearing.
\end{itemize}

\subsubsection{Example: Shadow-Cat (Lunara)}

\begin{itemize}
\item \textbf{Cap:} 3
\item \textbf{Body:} 4, \textbf{Wits:} 2, \textbf{Spirit:} 3, \textbf{Presence:} 1
\item \textbf{Harm:} 5
\item \textbf{Traits:} Silent Movement, Spirit Sense, Shadow Step (once per scene, teleport between shadows within Near range)
\end{itemize}

\subsection{Designing a Sentient Thiasos}

\subsubsection{Core Statistics}

A sentient Thiasos has no Harm track and cannot act independently. It has:

\begin{itemize}
\item \textbf{Personality:} A brief description of its voice, temperament, and quirks.
\item \textbf{Witness:} Always counts as a witness for Rites.
\item \textbf{Communication:} Can speak aloud or only to the Runekeeper.
\item \textbf{Abilities:} One to two special abilities (see below)
\end{itemize}

\subsubsection{Sample Abilities}

\begin{itemize}
\item \textbf{Perfect Recall:} Remembers everything it has witnessed. Can recall verbatim conversations from the past month.
\item \textbf{Warning:} Vibrates or hums when danger is near. The Runekeeper cannot be surprised.
\item \textbf{Advice:} Once per session, the Runekeeper may ask for advice on a roll. The GM may grant +1 die (good advice) or gain 1 SB (bad advice). The Thiasos chooses.
\item \textbf{Recording:} Can record a single truth or conversation and replay it later.
\item \textbf{Sensory Extension:} The Runekeeper can see or hear through the Thiasos within Near range.
\item \textbf{Impersonation:} Can speak in the voice of anyone whose blood it has tasted.
\end{itemize}

\subsubsection{Example: The Raven's Quill (Ikasha)}

\begin{itemize}
\item \textbf{Personality:} Secretive, amused, fond of riddles. Speaks in rustles.
\item \textbf{Witness:} Always counts as a witness.
\item \textbf{Communication:} Speaks only to the Runekeeper, in whispers.
\item \textbf{Abilities:} Perfect Recall, Recording (can record one secret and replay it later)
\end{itemize}

\subsection{Designing a Unified Thiasos (Codex and Thiasos as One)}

\subsubsection{Requirements}

\begin{itemize}
\item The object must be crafted or acquired through a downtime project (one week, Craft test DV 4).
\item The Runekeeper must perform a bonding ritual (one day, requires a threshold and a witness).
\item The Runekeeper marks +2 Obligation immediately.
\end{itemize}

\subsubsection{Core Statistics}

A unified Thiasos combines the traits of a Codex and a sentient Thiasos. It has:

\begin{itemize}
\item \textbf{Form:} Weapon, shield, staff, bell, or other object.
\item \textbf{Personality:} As sentient Thiasos.
\item \textbf{Witness:} Always counts as a witness.
\item \textbf{Codex:} Rites engraved on its surface. Cannot be destroyed by mundane means.
\item \textbf{Abilities:} One to two special abilities (may be drawn from living or sentient lists, but adjusted for object form)
\end{itemize}

\subsubsection{Example: Dawnblade (Oath of Flame \& Light)}

\begin{itemize}
\item \textbf{Form:} Longsword with Rites engraved along the fuller.
\item \textbf{Personality:} Stern, honorable, judgmental. Speaks in proverbs.
\item \textbf{Witness:} Always counts as a witness.
\item \textbf{Codex:} Rites of the Oath of Flame and Light.
\item \textbf{Abilities:} Warning (hums when oath-breakers are near), Advice (once per session, grant +1 die on a roll involving justice or honor)
\end{itemize}

\subsection{Balancing Custom Thiasos}

A Thiasos should not overshadow the Runekeeper. It is a companion, not a replacement.

\begin{itemize}
\item A living Thiasos should have fewer actions than the Runekeeper (once per scene independent action, otherwise requiring commands).
\item A sentient Thiasos should not have combat abilities. It is a witness, a counselor, a tool—not a warrior.
\item A unified Thiasos trades the safety of a separate Codex for the power of a bonded object. It is powerful, but its loss is catastrophic.
\end{itemize}

\subsubsection{Power Budget}

When designing a custom Thiasos, use the following budget:

\begin{itemize}
\item \textbf{Cap 2:} Two traits or one powerful trait.
\item \textbf{Cap 3:} Three traits or two powerful traits.
\item \textbf{Cap 4:} Four traits or three powerful traits.
\item \textbf{Cap 5:} Five traits or four powerful traits.
\end{itemize}

A "powerful trait" might be Regeneration, Terrifying Presence, or Perfect Recall. A standard trait might be Tracker, Silent Movement, or Warning.

\subsection{Ritual of Bonding}

Acquiring a new Thiasos—or bonding with a unified Codex-Thiasos—requires a ritual.

\subsubsection{Requirements}

\begin{itemize}
\item \textbf{Threshold:} A place significant to the Patron (shrine, crossroads, battlefield, grave).
\item \textbf{Witness:} A person who will observe the bonding.
\item \textbf{Offering:} Something of value to the Patron (blood, a memory, a sworn oath, a crafted object).
\item \textbf{Time:} One day (significant downtime).
\end{itemize}

\subsubsection{Procedure}

\begin{enumerate}
\item The Runekeeper and the prospective Thiasos stand at the threshold.
\item The witness observes.
\item The Runekeeper makes an offering and speaks the terms of the covenant.
\item The Runekeeper rolls Spirit + Lore (DV 4). On a success, the bond is formed. On a partial, the bond is formed but the Thiasos is Wary (suffers -1 die to its first action each scene) until the Runekeeper proves their worth. On a miss, the bond fails; the offering is consumed, and the Runekeeper marks 1 Obligation.
\end{enumerate}

\subsection{Sample Custom Thiasos}

\subsubsection{The Copper Hound (Mykkiel)}

\textbf{Form:} Living—a copper-coated mechanical hound with emerald eyes.

\textbf{Cap:} 3

\textbf{Statistics:} Body 4, Wits 3, Spirit 2, Presence 1, Harm 5

\textbf{Traits:}
\begin{itemize}
\item \textit{Tracker:} +2 dice to track a target by scent.
\item \textit{Oath-Sense:} Can detect whether a creature has broken a sworn oath within the past day.
\item \textit{Metallic Hide:} +1 Armor against mundane attacks.
\end{itemize}

\textbf{Personality:} Loyal, literal, slightly obsessive. It speaks in short sentences and repeats important commands.

\subsubsection{The Singing Lantern (The Pale Shepherd)}

\textbf{Form:} Sentient—a bronze lantern that burns with a cold, grey flame.

\textbf{Personality:} Melancholy, patient, soft-spoken. It hums when death is near.

\textbf{Abilities:}
\begin{itemize}
\item \textit{Warning:} Hum when someone within Near range is about to die (within the next hour).
\item \textit{Grey Light:} The lantern's light reveals invisible spirits and incorporeal creatures.
\item \textit{Final Witness:} Can record the last words of a dying creature and replay them once.
\end{itemize}

\subsubsection{The Storm-Drum (Xhak'Thul)}

\textbf{Form:} Unified—a hand-drum made from storm-felled wood and the hide of a beast killed by lightning.

\textbf{Personality:} Loud, impatient, energetic. It speaks in drumbeats that the Runekeeper can understand.

\textbf{Abilities:}
\begin{itemize}
\item \textit{Warning:} Beats faster when danger approaches.
\item \textit{Thunderclap:} Once per scene, the Runekeeper may strike the drum to create a deafening sound. Enemies within Near range test Resolve (DV 3) or lose their next minor action.
\item \textit{Storm-Sense:} The drum can predict the weather for the next hour with perfect accuracy.
\end{itemize}

\subsection{GM Guidance: Thiasos as Narrative Engine}

A Thiasos is not just a stat block. It is a character. Use it to:

\begin{itemize}
\item Deliver exposition or hints when the party is stuck.
\item Create conflict—the Thiasos may demand actions the Runekeeper is unwilling to take.
\item Foreshadow danger—the Thiasos can sense threats before they arrive.
\item Remember the past—the Thiasos has witnessed everything. It may reveal things the Runekeeper has forgotten—or wishes they had.
\end{itemize}

The thread held. The water carried. That is not a Rite. That is grace.
